\scanpage{8}
% L1
Soit $C'(\deleted{\unsure{p_0,p_1}}\unsure{\nu_0,\nu_1|v_0,v_1})$ la cat. discrète définie par les couples
% L2
$(\alpha_0,\alpha_1)$, \deleted{satisfaisant chacun des \unsure{satisfais|satisfait}} avec \deleted{\unsure{les}}
% L3
\[\text{a)}\qquad(\alpha_0+2)(\alpha_1\unsure{+|-}2)(\alpha_0+\alpha_1)\quad\text{inversible}\]
% L4
et chaque $\alpha_i$ satisfaisant les conditions suivantes
% L5
\[\text{b)}\qquad\Psi_{\nu_i}(\alpha_i)=0\qquad i=0,1\]
% L6
\deleted{\unsure{c)}}
% L7-8
\deleted{Supposons enfin que, si l'un des quand l'un des $\nu_i$ est tel que $\nu_i=2p$ (où $p$ nb premier) impair,}
% L9-11
\deleted{c) Si $\nu_i\cdot1_S$ \deleted{est de la forme $2p$ ($p$ nb premier) \unsure{alors on a} $p$} \inserted{$S$ est} \deleted{est inv. sur $S$} est inversible, \uline{ou} de la forme $2p$ ($p$ nb premier), \inserted{avec $2\cdot1_S$ inv. si $p$ impair,} (ou $\nu_i=p$ (nb premier impair)).}
% slanted
% L11a
\margin{\deleted{Condition ne dépendant que de $S$ et $\nu_0,\nu_1$ pas de l'objet particulier choisi.}}
% L12
Alors on a une équivalence de catégories
% L13
\[ C_{(\nu_0,\nu_1)}\longrightarrow C'(\nu_0,\nu_1). \]
\begin{center}\rule{.11\linewidth}{.4pt}\end{center}

% L14
En effet, il faut prouver, \deleted{\unsure{en supposant}}\inserted{si $\Gamma=\unsure{G|g|G_2|\mathcal G}$}, et
% L15
si $(G,\varphi)$ \inserted{$\in\mathrm{Ob}\ C$} et $(\alpha_0,\alpha_1)$ \inserted{$\in\mathrm{Ob}\ C'$} se correspondent, que
% L16
\deleted{A) les conditions a) b) c) sont satisfaites, ssi on}
% L17
l'on a la condition \deleted{a) ci-dessus sur $\rho_0$ et $\rho_1$}
% L18
\[\deleted{\text{ssi}}\quad\rho_i^{\deleted{\unsure{\rho_0}}\nu_i}=1\quad\text{et}\quad\forall s\in S,\quad\rho_{is}\ \text{d'ordre }\nu_i\text{ exactem\unsure{t}}\tag*{(*)}\]
% L19
ssi on a \unsure{'}
% L20
\[\Psi_{\nu_i}(\alpha_i)=0\quad\deleted{\text{et c) concernant }\nu_i}\tag*{(**)}\]

% L21
\margin{\uline{réf}}\deleted{Or il est} Or on sait que \inserted{(pour $\nu_i\geq3$)}, \deleted{suivant que $\nu_i$}
% L22-23
\deleted{est de la forme $2p$ ($p$ premier) ou $p$ ($p$ premier impair), ou ni l'un ni l'autre, que} \inserted{\deleted{premier, ou}} ça s'exprime
% L24
par la condition $\Psi_{\nu_i}(\alpha_i)=0$ sur l'invariant $\alpha_i$
% L25
de $\rho_i$, plus la condition $\nu_i1_S$ inv. si $\nu_i$ non premier
% L26
:
% L27
\deleted{Cas 1\textsuperscript{o}) $2$ inv. sur $S$ si $p$ impair / mais on a}
% L28
\deleted{Cas 2\textsuperscript{o}) aucune condition suppl. \qquad supp. ou $\nu_i$ inv.}
% L29
\deleted{Cas 3\textsuperscript{o}) $\nu_i$ inv. sur $S$ \qquad / ou 1\textsuperscript{o}) supprimé}
% L30
Ce qui prouve l'équivalence.
% archive sheet number: 7
